% ---- pg-416 （书页 404） ----

(a) 画因子得分图, 找可疑观测值.
计算对各观测值的标准化得分和 4.6 节所述的平方距离.
(b) 试用最大方差旋转.
2. 执行包括最大方差旋转在内的极大似然因子分析.
3. 比较从两个因子分析得到的解.
(a) 载荷的分组方式是否相同?
(b) 画出从主成分得到的因子得分图, 对比从极大似然分析得到的图形.
4. 对其他的公共因子数 $m$, 重复前三个步骤. 有必要将额外的因子放入来理解和阐述该数据吗?
5. 对大数据集, 将它们分为两半并且对每一半都执行因子分析. 比较这两个结果, 并与从完整数据集所得结果做比较, 以检查解的稳定性. (数据剖分时可以把前一半放在一组中, 把后一半放在另一组中. 这能揭示随时间而发生的变化.)

\noindent\textbf{例 9.14}\quad (鸡的骨骼数据的因子分析)

我们介绍对白来亨鸡的骨骼和颅骨测量的几个因子分析的结果, 原始数据取自邓恩(Dunn, 见文献[5]). 邓恩数据的因子分析, 起初由赖特(Wright, 见文献[15])研究, 他从跟我们不同的相关矩阵着手分析.

全数据集由对鸡骨骼大小的 $n=276$ 个测量值(数据)组成:
\[
\begin{array}{r@{\ }l}
\text{头：} & \left\{\begin{array}{l}X_1=\text{颅骨长度}\\ X_2=\text{颅骨幅宽}\end{array}\right.\\[10pt]
\text{腿：} & \left\{\begin{array}{l}X_3=\text{股骨长}\\ X_4=\text{胫骨长}\end{array}\right.\\[10pt]
\text{翅：} & \left\{\begin{array}{l}X_5=\text{肱骨长}\\ X_6=\text{尺骨长}\end{array}\right.
\end{array}
\]

把样本相关矩阵
\[
\mathbf{R}=\begin{bmatrix}
1.000&0.505&0.569&0.602&0.621&0.603\\
0.505&1.000&0.422&0.467&0.482&0.450\\
0.569&0.422&1.000&0.926&0.877&0.878\\
0.602&0.467&0.926&1.000&0.874&0.894\\
0.621&0.482&0.877&0.874&1.000&0.937\\
0.603&0.450&0.878&0.894&0.937&1.000
\end{bmatrix}
\]

用主成分法和极大似然法对 $m=3$ 因子模型作因子分析. 结果显示于\hyperref[tab:9.10]{表 9.10}.\footnote{{\footnotesize 注意在极大似然解中胫骨长的估计特殊方差为 0.00. 这表明极大化似然函数可以产生海伍德情况. 重做我们的结果的读者若使用 SAS 或类似软件, 应该试一下 Hey(wood) 选择.}}

旋转之后, 求解的两种方法出现略为不同的结果. 将注意力集中于主成分法和被解释的样本总方差的累积比, 我们看到三因子解是有道理的. 第三个因子解释附加样本变差的一个“显著的”比例. 第一个因子表现为体格大小因子, 为翅和腿的大小所左右. 第二和第三因子, 总体上描绘颅骨大小, 可以分别给予与变量相同的名字, 即颅骨幅宽和颅骨长度.
